\chapter{Marco teórico}

\section{Dinámica de Fluidos Computacional}

\subsection{Método de los volúmenes Finitos}

Acá mostramos las ecuaciones de Navier Stokes y empezamos a tirar integrales hasta llegar al MVF.

\subsection{Resolución de la turbulencia}

Acá vamos a desarrollar la teoría detrás de RANS y LES

\begin{itemize}
	\item RANS - kEpsilon standard
	\item RANS - kEpsilon realizable
	\item RANS - kEpsilon Sogachev
	\item LES - Smagorinsky
	\item LES -  oneEquation eddy viscosity
\end{itemize}

\subsection{OpenFOAM}

Acá vamos a hablar brevemente de la estructura de datos de OpenFOAM. En particular de como agregar términos fuentes.

\section{Capa Límite Atmosférica y Estratificación Térmica}

Acá vamos a describir toda la data de la capa límite atmosférica. 

\subsection{Capa Límite Atmosférica}

Comenzamos con descripciones generales.
Luego tenemos agregamos perfiles de viento típicos. 

\subsection{Efecto de la Temperatura}

Aca viene la estratificación

\subsection{Capa Superficial Atmosférica vs. Capa Límite Atmosférica}

Los 2 modelos posibles

\subsection{Capa Límite Atmosférica Convencionalmente Neutra (CNABL)}

\section{Simulación de la CNABL}

Acá hay que mencionar una serie de cosas:

\subsection{Hipótesis de Bousinesqq para plantear las ecuaciones con temperatura}
\subsection{Generación de Condiciones de borde reales}
Precursor vs tendencias desde WRF
\subsection{En casos de RANS incorporación de Temperatura vs limitación de la longitud de mezcla}
estacionario vs transitorio




\section{Aerodinámica de aerogeneradores}

El modelo de actuador discal fue propuesto originalmente por Rankine \citep{rankine1865} y Froude \citep{froude1889}

\subsection{Teoría Unidimensional de Momento}

\subsection{Teoría BEM}

\subsection{Aerodinámica real de aerogeneradores}

\subsection{Efecto estela en aerogeneradores}

\subsection{Interacción con la CNABL}

\section{Aerodinámica de aerogeneradores flotantes}

\subsection{Consecuencias del movimiento de la plataforma}

\subsection{Efecto estela en aerogeneradores flotantes}

\section{Modelos simplificados de aerogeneradores}

Acá hacemos una descripción de los modelos existentes, full rotor, actuador sectorial, actuador de superficie
etc. Luego vamos de lleno con los que nos interesan. 

AD

AL

Tip factor

Root Factor

gaussiana



\section{Técnicas de Análisis de Resultados}

\subsection{Descomposición Ortogonal Propia (POD)}
\subsection{Promediado en fase}
